% © 2026 Ing. David Jaroš — CC BY-NC-ND 4.0
%
% This work is licensed under a Creative Commons Attribution-NonCommercial-NoDerivatives
% 4.0 International License (CC BY-NC-ND 4.0).
%
% License History: Earlier drafts (up to v0.3) were released under CC BY 4.0.
% From v0.4 onward, all material is released under CC BY-NC-ND 4.0 to protect
% the integrity of the theoretical work during ongoing academic development.
%
% See LICENSE.md for full license text.
%
% Track: EW — Electroweak sector
% File: research_tracks/EW/higgs_from_theta.tex
% Purpose: First-step analysis of spontaneous symmetry breaking in UBT via S[Θ].
%          Checks whether Coleman-Weinberg mechanism gives VEV ⟨Θ₀⟩ ≠ 0.
% Status: OPEN — canonical action checked; CW computation needed.
% Date: 2026-05-17

\documentclass[12pt,a4paper]{article}
\usepackage[a4paper,margin=2.5cm]{geometry}
\usepackage{amsmath,amssymb,amsthm,booktabs}
\usepackage{tcolorbox}
\tcbuselibrary{skins,breakable}
\usepackage{hyperref}

\newtcolorbox{statusbox}[1][]{%
  colback=gray!8!white,colframe=black!60,arc=3pt,boxrule=1pt,
  fonttitle=\bfseries,title=Status Summary,#1}
\newtcolorbox{openbox}[1][]{%
  colback=yellow!5!white,colframe=orange!55!black,arc=3pt,boxrule=1pt,
  fonttitle=\bfseries,title=Open Problem,#1}

\newtheorem{theorem}{Theorem}[section]
\newtheorem{lemma}[theorem]{Lemma}
\newtheorem{proposition}[theorem]{Proposition}

\title{Higgs Mechanism from $S[\Theta]$:\\[4pt]
       First Step — Coleman-Weinberg Analysis}
\author{Ing.\ David Jaro\v{s}}
\date{2026-05-17}

\begin{document}
\maketitle

\begin{abstract}
We investigate whether the UBT action $S[\Theta]$ admits spontaneous symmetry
breaking (SSB) via the Coleman--Weinberg (CW) mechanism.  The first step is
to determine whether $S[\Theta]$ contains an explicit scalar potential $V(\Theta)$.
We find that the canonical action has no tree-level Higgs-type potential;
SSB must therefore arise from radiative corrections.  The first-principles
derivation of the CW effective potential for $\Theta$ is identified as
\textbf{[OPEN]} (Gap EW-2), with the first non-trivial computation step described.
$W/Z$ masses from SSB remain \textbf{[OPEN]}.
\end{abstract}

\tableofcontents

%──────────────────────────────────────────────────────────────────────────────
\section{Tree-Level Action Check}
\label{sec:tree_level}
%──────────────────────────────────────────────────────────────────────────────

The canonical UBT action is (from \texttt{canonical/UBT\_canonical\_main.tex}):
\[
  S[\Theta] = \int d^4x\,d\psi\;
  \bigl[\tfrac{1}{2}\,(\nabla^\dagger\nabla\Theta)^2
  + \mathcal{L}_{\rm int}\bigr],
\]
where $\mathcal{L}_{\rm int}$ describes gauge interactions.

\begin{proposition}[Tree-level potential check]
\label{prop:no_tree_potential}
The canonical action $S[\Theta]$ does \emph{not} contain an explicit
quartic or mass potential $V(\Theta)$ of the Higgs-type
$V = -\mu^2|\Theta|^2 + \lambda|\Theta|^4$.
\end{proposition}

\begin{proof}
Direct inspection of $S[\Theta]$: the kinetic term
$(\nabla^\dagger\nabla\Theta)^2$ produces only derivative couplings.
The interaction term $\mathcal{L}_{\rm int}$ describes minimal gauge coupling
but no non-derivative self-interaction of $\Theta$.
There is no free parameter for a mass $\mu^2$ or quartic coupling $\lambda$.
\end{proof}

\begin{openbox}
If $S[\Theta]$ has no tree-level potential, SSB must arise from:
(a) radiative corrections (Coleman--Weinberg mechanism), or
(b) boundary conditions on $S^1_\psi$ (Hosotani mechanism), or
(c) a non-perturbative condensate.
All three routes are open; (a) is pursued in this note.
\end{openbox}

%──────────────────────────────────────────────────────────────────────────────
\section{Coleman-Weinberg Effective Potential}
\label{sec:coleman_weinberg}
%──────────────────────────────────────────────────────────────────────────────

\subsection{Setup}

The Coleman--Weinberg effective potential for a field $\phi$ with
field-dependent mass $M(\phi)$ is
\begin{equation}
  \label{eq:CW_potential}
  V_{\rm CW}(\Theta) = \frac{1}{64\pi^2}\,
  \mathrm{Str}\bigl[M^4(\Theta)\bigl(\ln\tfrac{M^2(\Theta)}{\mu^2}
  - \tfrac{3}{2}\bigr)\bigr],
\end{equation}
where $\mathrm{Str}$ denotes the supertrace over all modes.

For UBT, the field-dependent mass of KK mode $n$ is
\[
  M_n(\Theta) = \frac{n}{R_\psi}\,f(\Theta),
\]
where $f(\Theta)$ encodes the coupling of $\Theta$ to the KK mass.

\subsection{Condition for SSB}

Spontaneous symmetry breaking occurs if
\[
  \frac{\partial V_{\rm CW}}{\partial |\Theta|}\bigg|_{\Theta=0} < 0,
\]
i.e.\ if the origin $\Theta = 0$ is a local maximum of $V_{\rm eff}$.

Expanding around $\Theta = \Theta_0 + \delta\Theta$ with $\Theta_0$ a background:
\begin{align}
  V_{\rm CW}(\Theta_0) &\sim
  \sum_n \frac{M_n^4(\Theta_0)}{64\pi^2}
  \Bigl(\ln\frac{M_n^2(\Theta_0)}{\mu^2} - \frac{3}{2}\Bigr).
\end{align}

For a VEV to exist, the one-loop mass-squared matrix must have a tachyonic
direction, i.e.\ the second derivative of $V_{\rm CW}$ with respect to
fluctuations of $|\Theta|$ must be negative at $|\Theta|=0$.

\subsection{First computation step}

The leading contribution at order $R_\psi^{-4}$ is
\[
  V_{\rm CW}^{(1)}(\Theta_0) \sim
  \frac{N_{\rm eff}}{64\pi^2}\,
  \frac{f^4(\Theta_0)}{R_\psi^4}
  \Bigl(\ln\frac{f^2(\Theta_0)}{R_\psi^2 \mu^2} - \frac{3}{2}\Bigr)
  \sum_n n^4\,e^{-n/\Lambda R_\psi},
\]
where $\Lambda$ is a UV regulator.
The sum $\sum_n n^4 e^{-n\epsilon}$ is regulated by the Epstein zeta function
(as in \texttt{research\_tracks/T3\_ALPHA/mellin\_insertion\_B.tex}).

\begin{openbox}
\textbf{Open step}: Compute $\partial^2 V_{\rm CW}/\partial|\Theta|^2|_{\Theta=0}$
from the UBT action.  The sign determines whether SSB occurs.
Current status: $f(\Theta)$ not derived from $S[\Theta]$; coupling between
$\Theta$ background and KK mass not computed.
This is Gap EW-2.
\end{openbox}

%──────────────────────────────────────────────────────────────────────────────
\section{Hosotani Mechanism: Alternative Route}
\label{sec:hosotani}
%──────────────────────────────────────────────────────────────────────────────

An alternative SSB route is the Hosotani (Wilson-line) mechanism:
boundary conditions on $S^1_\psi$ for gauge fields $A_\psi$ generate
an effective potential whose minimum breaks the gauge symmetry.

From \texttt{research\_tracks/EW/rpsi\_from\_action.tex}~[L1 conditional],
the Scherk--Schwarz projection gives $\lambda_{\rm eff} = g^2/2$.
Whether the Hosotani potential minimum coincides with the electroweak
VEV $v = 246\,\mathrm{GeV}$ is~\textbf{[OPEN]}.

%──────────────────────────────────────────────────────────────────────────────
\section{W/Z Masses: Conditional Chain}
\label{sec:wz_masses}
%──────────────────────────────────────────────────────────────────────────────

Assuming SSB with VEV $\langle\Theta_0\rangle = v\,\mathbf{1}_2$, the
$W$ and $Z$ masses arise from the gauge kinetic term:
\[
  \mathcal{L}_{\rm gauge} \supset
  g^2 v^2 W^+_\mu W^{-\mu} + \tfrac{1}{2}(g^2+g'^2)v^2 Z_\mu Z^\mu.
\]
This gives the standard Higgs mechanism results
$M_W = gv/2$ and $M_Z = \sqrt{g^2+g'^2}\,v/2$.

However, deriving $v$ from $S[\Theta]$ requires:
\begin{enumerate}
\item A first-principles CW potential computation (Gap EW-2);
\item The coupling $f(\Theta)$ determined from $S[\Theta]$;
\item The minimum of $V_{\rm CW}$ to lie at $|\Theta| = v \neq 0$.
\end{enumerate}
All three steps are~\textbf{[OPEN]}.

%──────────────────────────────────────────────────────────────────────────────
\begin{statusbox}
$S[\Theta]$ explicit tree-level $V(\Theta)$: \textbf{NO} (Prop.~\ref{prop:no_tree_potential})\\
Coleman-Weinberg SSB possible: \textbf{[OPEN]} — $f(\Theta)$ not derived\\
VEV $\langle\Theta_0\rangle \neq 0$: \textbf{[OPEN]} — conditional on CW computation\\
$W/Z$ masses from SSB: \textbf{[OPEN]} — conditional on VEV\\
Hosotani mechanism: \textbf{[OPEN]} — Hosotani potential minimum not computed\\
Higgs mass 125\,GeV: \textbf{[OPEN]} (not addressed here)\\
Gap EW-2 status: \textbf{OPEN} — first step identified
\end{statusbox}
%──────────────────────────────────────────────────────────────────────────────

\end{document}
